\documentclass[11pt]{article}

\usepackage[margin=1in]{geometry}
\usepackage{amsmath,amssymb,amsthm}
\usepackage{graphicx}
\usepackage{microtype}
\usepackage[hidelinks]{hyperref}

\newtheorem{theorem}{Theorem}
\newtheorem{lemma}[theorem]{Lemma}
\newcommand{\weak}{\sigma(X,X^*)}
\newcommand{\diam}{\operatorname{diam}}

\title{Weak Banach-Space Groups Are Fattening\
Exactly in Finite Dimension\\
\large A full candidate answer to the weak-topology branch of Question 1 in arXiv:2104.02695}
\author{Automated mathematics run \texttt{fa\_banach\_001}, lane 5}
\date{9 August 2026}

\begin{document}
\maketitle

\begin{center}
\fbox{\parbox{0.91\textwidth}{\textbf{Status.}
Candidate full solution to the explicit weak-topology branch of Open
Question 1, likely valid and awaiting human review.  The broader parts of
Question 1 concerning arbitrary Hausdorff or Polish groups are not claimed.}}
\end{center}

\section{Source question and definitions}

The source is D.~Ravasini, \emph{Compactivorous Sets in Banach Spaces},
arXiv:2104.02695, later published in \emph{Proceedings of the American
Mathematical Society} \textbf{150} (2022), 2121--2129
\cite{Ravasini}.  On page 8, Open Question 1 asks in particular what can be
said about weak topologies of Banach spaces without the Schur property.
The entire open-question section is reproduced below; the present packet
answers precisely the weak-topology clause of item (1).

\begin{figure}[ht]
\centering
\includegraphics[width=\textwidth]{figures/open_problem_crop.png}
\caption{Open Questions 1 and 2 in the source paper, page 8.}
\end{figure}

For a Hausdorff topological group $G$, a set $E\subseteq G$ is
\emph{compactivorous} if, for every compact $K\subseteq G$, there are an
open set $V\subseteq G$ and $g,h\in G$ such that
\[
 K\cap V\ne\varnothing,
 \qquad g(K\cap V)h\subseteq E.
\]
It is \emph{strongly compactivorous} if there is a closed neighbourhood $U$
of the identity such that, for every compact $K\subseteq U$, some $g,h\in G$
satisfy $gKh\subseteq E$.  The group is \emph{fattening} if every
compactivorous subset is strongly compactivorous.  For the additive group of
a Banach space, the two-sided translate $gAh$ is simply one translate of
$A$.

The source proves that an infinite-dimensional Banach space with the Schur
property is nonfattening in its weak topology.  It uses the Schur property to
turn weakly compact sets into norm-compact sets.  The theorem below removes
that hypothesis completely.

\section{Proof intuition}

There are two independent scales in a weakly topologised Banach space.
Weakly compact sets have small pieces in the norm metric: the classical
norm-fragmentation theorem says that every weakly compact set has a nonempty
relatively weak-open subset of arbitrarily small norm diameter.  Translating
one point of such a piece to the origin puts the whole piece in the norm unit
ball.  Thus the norm unit ball is compactivorous even though it has empty weak
interior in the typical infinite-dimensional case.

On the other hand, a basic weak neighbourhood restricts only finitely many
functionals.  Their common kernel is an infinite-dimensional direction on
which the neighbourhood is unbounded.  It therefore contains a two-point
compact set of diameter larger than the unit ball's diameter.  No translation
can fit that set into the unit ball, so the unit ball is not strongly
compactivorous.  Finite-dimensional spaces fall on the opposite side because
weak and norm topologies coincide and the resulting group is locally compact.

\section{The fragmentation input}

A metric $d$ \emph{fragments} a topological space $T$ if, for every nonempty
$A\subseteq T$ and every $\varepsilon>0$, there is an open $O\subseteq T$
such that $A\cap O\ne\varnothing$ and
$\diam_d(A\cap O)<\varepsilon$.

\begin{lemma}[Weak-compact norm fragmentation]\label{lem:fragment}
Let $X$ be a Banach space and let $K\subseteq X$ be weakly compact.  For every
$\varepsilon>0$ there is a weakly open set $V\subseteq X$ such that
\[
 K\cap V\ne\varnothing,
 \qquad \diam_{\|\cdot\|}(K\cap V)<\varepsilon.
\]
\end{lemma}

\noindent
This is the standard theorem that weakly compact subsets of Banach spaces are
fragmented by the norm; see Jayne--Namioka--Rogers
\cite{JayneNamiokaRogers}.  Their theorem has no separability or Schur
hypothesis.  Applying the definition of fragmentation to $A=K$ gives the
displayed formulation.  This is the only non-elementary external result used
below.

\section{Main theorem}

\begin{theorem}\label{thm:main}
Let $X$ be a real or complex Banach space, and give its additive group the weak
topology $\weak$.  Then the following are equivalent:
\begin{enumerate}
\item $(X,\weak)$ is fattening;
\item $X$ is finite-dimensional.
\end{enumerate}
More precisely, whenever $X$ is infinite-dimensional, its closed norm unit
ball $B_X$ is compactivorous in $(X,\weak)$ but is not strongly
compactivorous.
\end{theorem}

\begin{proof}
Suppose first that $X$ is infinite-dimensional.  Let $K\subseteq X$ be any
nonempty weakly compact set.  Apply Lemma~\ref{lem:fragment} with
$\varepsilon=1$ to obtain a weakly open $V\subseteq X$ such that
$K\cap V$ is nonempty and has norm diameter less than $1$.  Choose
$k_0\in K\cap V$.  Then
\[
 -k_0+(K\cap V)\subseteq B_X,
\]
because $\|k-k_0\|<1$ for every $k\in K\cap V$.  Since $K$ was arbitrary,
$B_X$ is compactivorous in the weak topology.

We next show that $B_X$ is not strongly compactivorous.  Let $U$ be an
arbitrary weak neighbourhood of $0$; the argument therefore also applies when
$U$ is closed.  Some basic weak neighbourhood
\[
 W=\{x\in X: |x_j^*(x)|<\delta\text{ for }j=1,\ldots,n\}
\]
is contained in $U$, for suitable $x_1^*,\ldots,x_n^*\in X^*$ and
$\delta>0$.  The subspace
\[
 Y=\bigcap_{j=1}^n\ker x_j^*
\]
has finite codimension and is nonzero (indeed infinite-dimensional), so
$Y\subseteq W$ contains a vector $z$ with $\|z\|>1$.  The two-point set
$K=\{-z,z\}$ is weakly compact and lies in $U$.  Every translate of $K$ has
diameter $2\|z\|>2$, whereas every subset of $B_X$ has diameter at most $2$.
Consequently no translate of $K$ lies in $B_X$.  Since this happens inside
every weak neighbourhood $U$ of $0$, the set $B_X$ is not strongly
compactivorous.  Hence $(X,\weak)$ is not fattening.

Finally, suppose $X$ is finite-dimensional.  The weak and norm topologies
coincide.  For completeness, let $E\subseteq X$ be compactivorous and apply
the definition to a compact closed ball $C$ with nonempty interior.  There
are an open $V$ and a translate $t+(C\cap V)\subseteq E$, with
$C\cap V\ne\varnothing$.  Since $\operatorname{int}C$ is dense in $C$, the
set $V\cap\operatorname{int}C$ is nonempty.  Thus $C\cap V$ contains a
translate $c+W$ of an open neighbourhood $W$ of $0$, and $E$ contains
$t+c+W$.  Choose a closed neighbourhood $U\subseteq W$.  For every compact
$K\subseteq U$, the fixed translate $t+c+K$ lies in $E$.  Hence $E$ is
strongly compactivorous, so $X$ is fattening.  (This is the elementary
finite-dimensional instance of the source paper's observation that locally
compact groups are fattening.)
\end{proof}

\section{Verification notes}

The proof uses exactly one imported geometric fact, Lemma~\ref{lem:fragment}.
Its quantifiers match the compactivorous definition: the relatively weak-open
piece is the intersection of $K$ with an ambient weakly open set, and the
theorem applies to arbitrary weakly compact $K$, without metrizability or
separability.  The remainder is elementary.  Translation preserves diameter;
a basic weak neighbourhood contains the common kernel of finitely many
functionals; and that kernel is nonzero in infinite dimension.  The witness
$\{-z,z\}$ is finite, hence compact in every Hausdorff topology on $X$.

No numerical or symbolic experiment is part of the proof.  A separate
verification report in the packet records a dependency and definition audit.
The main human-review target is the precise norm-fragmentation theorem and its
application to the source paper's group-theoretic definition.

\section{Scope and novelty check}

This is a full answer to the explicit weak-topology branch of Open Question 1:
there is no additional infinite-dimensional class to find, because every
infinite-dimensional weak Banach-space group is nonfattening.  It does not
answer the other clauses about arbitrary Hausdorff groups or Polish groups,
and it does not address Open Question 2 about Haar-null closed convex sets.

The bounded novelty search was performed on 9 August 2026.  The run indexes
were searched for arXiv:2104.02695, \texttt{compactivorous},
\texttt{fattening}, and the weak-topology/Schur terminology.  Web and arXiv
searches used the exact source title and the phrases
``weak topologies of Banach spaces'', ``fattening'',
``without the Schur property'', and close combinations with
``compactivorous''.  The searches surfaced the source article, its journal
version, and later work on Haar-null convex sets, but no later resolution of
this weak-topology question.  This is evidence only, not a claim of exhaustive
bibliographic novelty; specialist review remains necessary.

\section{Human-review recommendation}

Review as a focused full solution to the weak-topology subquestion.  First
confirm that the cited fragmentation theorem is understood in its standard
nonseparable form.  Then check the two short definition bridges: translating
a small-diameter relatively weak-open piece into $B_X$, and using a two-point
set in the common kernel to refute every possible strong witness
neighbourhood.  If those checks pass, the argument yields the exact
finite-dimensional characterization claimed in Theorem~\ref{thm:main}.

\begin{thebibliography}{9}

\bibitem{Ravasini}
D. Ravasini,
\emph{Compactivorous Sets in Banach Spaces},
Proc. Amer. Math. Soc. \textbf{150} (2022), 2121--2129;
arXiv:2104.02695,
\href{https://doi.org/10.1090/proc/15851}{doi:10.1090/proc/15851}.

\bibitem{JayneNamiokaRogers}
J. E. Jayne, I. Namioka, and C. A. Rogers,
\emph{Topological Properties of Banach Spaces},
Proc. London Math. Soc. (3) \textbf{66} (1993), 651--672,
\href{https://doi.org/10.1112/plms/s3-66.3.651}{doi:10.1112/plms/s3-66.3.651}.

\end{thebibliography}

\end{document}
